% StoHyMoLAP / RAMIS paper template
% Suggested target style: Elsevier elsarticle / Results in Engineering-like structure.
% Compile with: latexmk -pdf paper_stohymolap_template.tex

\documentclass[preprint,12pt]{elsarticle}

\usepackage[utf8]{inputenc}
\usepackage[T1]{fontenc}
\usepackage{amsmath,amssymb}
\usepackage{graphicx}
\usepackage{booktabs}
\usepackage{array}
\usepackage{multirow}
\usepackage{siunitx}
\usepackage{xcolor}
\usepackage{hyperref}
\usepackage{lineno}

\journal{Results in Engineering}

% -----------------------------------------------------------------------------
% Macros
% -----------------------------------------------------------------------------
\newcommand{\Qobs}{Q_{\mathrm{obs}}}
\newcommand{\Qsim}{Q_{\mathrm{sim}}}
\newcommand{\NSE}{\mathrm{NSE}}
\newcommand{\KGE}{\mathrm{KGE}}
\newcommand{\RMSE}{\mathrm{RMSE}}
\newcommand{\MAE}{\mathrm{MAE}}
\newcommand{\PBIAS}{\mathrm{PBIAS}}
\newcommand{\warningbackend}{\textbf{Backend note:} E7/E8 must not be described as real recurrent GRU models unless TensorFlow/PyTorch evidence is regenerated. In the current artefacts they use a fallback MLP on flattened sequences.}

\begin{document}

\begin{frontmatter}

\title{Hybrid stochastic hydrological modelling with hydrological memory and machine learning for daily rainfall--runoff simulation}

\author[inst1]{Author One}
\author[inst2]{Author Two}
\address[inst1]{Institution, City, Country}
\address[inst2]{Institution, City, Country}

\begin{abstract}
% Draft after Results/Discussion are stable.
% Suggested content:
% 1) Problem: daily discharge simulation remains difficult under stochastic rainfall--runoff variability.
% 2) Aim: evaluate RAMIS/StoHyMoLAP-inspired deterministic, stochastic, ML, and hybrid configurations.
% 3) Methods: E0--E8 ablation matrix, common validation window, anti-leakage audit, metrics NSE/KGE/RMSE/MAE/PBIAS/R2.
% 4) Results: best common-window model and main ablation effects.
% 5) Limitation: GRU-labelled models currently executed with fallback MLP backend; stochastic bands show limited dispersion if applicable.
% 6) Implication: hybrid physical--data-driven inputs can improve point discharge prediction when carefully audited.
\end{abstract}

\begin{keyword}
Stochastic hydrological model \sep HyMoLAP \sep rainfall--runoff modelling \sep hybrid models \sep hydrological memory \sep machine learning \sep uncertainty
\end{keyword}

\end{frontmatter}

\linenumbers

% -----------------------------------------------------------------------------
\section{Introduction}
% Suggested structure:
% - Importance of accurate daily streamflow simulation.
% - Limitations of purely deterministic hydrological models under random fluctuations.
% - Limitations of standalone ML when physical memory/process information is weak.
% - Hybridization context: using hydrological model outputs/statistical properties as ML inputs.
% - Gap: reproducible ablation of RAMIS/StoHyMoLAP-inspired physical memory, Lévy stochasticity, and ML/hybrid post-processing.
% - Contributions:
%   (i) common-window E0--E8 experimental matrix;
%   (ii) baseflow/memory and stochastic ablation;
%   (iii) hybrid XGB/GRU-labelled quantile configurations;
%   (iv) anti-leakage and backend audit;
%   (v) paper-ready figures and reproducibility package.

% -----------------------------------------------------------------------------
\section{Materials and methods}

\subsection{Study area and hydro-meteorological data}
% Describe basin, period, forcing variables, observed discharge, temporal resolution.
% Include train/validation split and missing-data treatment.

\subsection{RAMIS / HyMoLAP-inspired deterministic rainfall--runoff formulation}
% Present deterministic core and effective precipitation logic.
% Keep equations consistent with the implemented model.

\subsection{Hydrological memory and baseflow reservoir}
% Explain the added baseflow/memory component.
% Link to E0 vs E1 and E2 vs E3 comparisons.

\subsection{Lévy-driven stochastic extension}
% Explain stochastic perturbation and intended uncertainty representation.
% State how stochastic summaries enter the experimental matrix.
% Avoid claiming well-calibrated uncertainty bands if manifest reports subdispersion.

\subsection{Machine learning and hybrid configurations}
% Describe E4 pure ML, E5/E6 hybrid XGB, E7/E8 GRU-labelled hybrids.
% Include backend note explicitly.
\warningbackend

\subsection{Experimental matrix}
\begin{table}[htbp]
\centering
\caption{Experimental matrix used for the common-window comparison. Replace values with the final table exported from the reproducibility package.}
\label{tab:experimental_matrix}
\begin{tabular}{llll}
\toprule
Experiment & Family & Main input & Purpose \\
\midrule
E0 & Physical & deterministic, no baseflow & baseline physical response \\
E1 & Physical & deterministic with baseflow & hydrological memory effect \\
E2 & Stochastic & Lévy, no baseflow & stochastic effect without memory \\
E3 & Stochastic & Lévy with baseflow & stochastic + memory effect \\
E4 & ML & pure XGB & data-driven baseline \\
E5 & Hybrid & mean physical/stochastic features + XGB & hybrid mean-feature effect \\
E6 & Hybrid & quantile features + XGB & hybrid uncertainty-feature effect \\
E7 & Hybrid & mean features + GRU-labelled fallback & sequence-labelled hybrid mean effect \\
E8 & Hybrid & quantile features + GRU-labelled fallback & sequence-labelled hybrid quantile effect \\
\bottomrule
\end{tabular}
\end{table}

\subsection{Evaluation metrics and common validation window}
% Define NSE, KGE, RMSE, MAE, PBIAS, R2 if used.
% State that all comparisons use exact common intersection and missing Qobs are excluded.

\subsection{Leakage and reproducibility checks}
% Report anti-leakage audit status and interpretation.
% Mention tests passed and generated manifests.

% -----------------------------------------------------------------------------
\section{Results}

\subsection{Common-window model ranking}
% Use table_final_ranking_common_window.csv.
% Main claim: identify best model by NSE/KGE over n_eval common samples.

\begin{table}[htbp]
\centering
\caption{Final ranking over the common validation window. Convert the CSV table to LaTeX before submission.}
\label{tab:final_ranking_common}
\begin{tabular}{lrrrrrrl}
\toprule
Model & NSE & KGE & RMSE & MAE & PBIAS & R$^2$ & Backend \\
\midrule
% Fill from outputs/reports/phase34/table_final_ranking_common_window.csv
\bottomrule
\end{tabular}
\end{table}

\subsection{Ablation effects of baseflow, stochasticity and hybridization}
% Use table_ablation_effects_common_window.csv.
% Discuss deltas in NSE/KGE/RMSE/MAE/PBIAS.

\subsection{Hydrograph behaviour in validation}
\begin{figure}[htbp]
\centering
\includegraphics[width=\linewidth]{outputs/figures/phase34/hydrograph_validation_common.png}
\caption{Observed and simulated validation hydrographs for selected physical, ML and hybrid configurations over the common validation window.}
\label{fig:hydrograph_common}
\end{figure}

\subsection{Flow-duration and scatter diagnostics}
\begin{figure}[htbp]
\centering
\includegraphics[width=0.95\linewidth]{outputs/figures/phase34/fdc_common.png}
\caption{Flow-duration curves comparing observed and simulated discharges for the selected configurations.}
\label{fig:fdc_common}
\end{figure}

% Add selected scatter figures or a composite if later generated.

\subsection{Error structure by hydrological regime}
% Use residuals_by_regime_*.png and table_regime_metrics_selected.csv.
% Interpret low/medium/high flow behaviour.

\subsection{Stochastic reference and uncertainty-band limitations}
% Use stochastic_reference_E2_E3.png and uncertainty table.
% State fallback/subdispersion if reported by manifest.

% -----------------------------------------------------------------------------
\section{Discussion}

\subsection{Why hydrological memory improves daily discharge simulation}
% Interpret baseflow/reservoir as slow-flow persistence and catchment memory.
% Connect to deterministic_baseflow and stochastic_baseflow ablations.

\subsection{Role of stochastic Lévy perturbations}
% Discuss whether stochasticity improves point metrics or mainly provides distributional features.
% Avoid claiming uncertainty calibration unless coverage/bands support it.

\subsection{Hybrid physical--machine learning gains}
% Discuss ML as correction/mapping from physical/stochastic features to Qsim.
% Compare pure ML vs hybrid models.

\subsection{Quantile-informed features and regime-dependent performance}
% Discuss why distributional summaries can help high/low flows.
% Compare E6 and E8 carefully.

\subsection{Comparison with previous stochastic HyMoLAP hybridization studies}
% Compare conceptually with Houenafa et al. (2025): stochastic model outputs, mean/quantile features, physical + ML hybridization.
% Emphasize differences: RAMIS/baseflow, common-window ablation, fallback backend caveat, current dataset.

\subsection{Limitations}
% Suggested points:
% - E7/E8 backend fallback is not a true recurrent GRU implementation in current environment.
% - Component plots skipped if Qfast/Qbase/Qtotal columns are absent.
% - Stochastic bands may show subdispersion or unavailable quantile columns.
% - Results depend on the common validation period and selected basin/data quality.
% - WARN audits require transparent reporting even if no FAIL is found.

\subsection{Implications and future work}
% Suggested points:
% - Re-run E7/E8 with TensorFlow/PyTorch GRU backend.
% - Export physical components Qfast/Qbase/Qtotal.
% - Improve stochastic coverage and calibration of quantile bands.
% - Test transferability across basins and hydroclimatic regimes.

% -----------------------------------------------------------------------------
\section{Conclusions}
% 4--6 concise conclusions:
% - Best common-window model and magnitude of gain.
% - Baseflow/memory contribution.
% - Quantile-informed hybrid contribution.
% - Backend caveat.
% - Reproducibility and future work.

\section*{Code and data availability}
% Provide repository URL, commit hash, scripts, outputs and manifests.

\section*{Declaration of competing interest}
The authors declare that they have no known competing financial interests or personal relationships that could have appeared to influence the work reported in this paper.

\section*{Acknowledgements}
% Optional.

\bibliographystyle{elsarticle-num}
\bibliography{references}

\end{document}
